\textbf{Phần 1. Trắc nghiệm lựa chọn}
\setcounterpageref{ex}{0}

\begin{ex} %[3X2V28]
Một cửa hàng buổi sáng bán được 24 gói bánh. Buổi chiều bán được nhiều hơn buổi sáng 12 gói bánh. Hỏi cả hai buổi cửa hàng bán được bao nhiêu gói bánh?
\choice
{$48$ gói bánh}
{\True $60$ gói bánh}
{$54$ gói bánh}
{$64$ gói bánh}
\loigiai{
Buổi chiều bán được: $24 + 12 = 36$ gói.\
Cả hai buổi bán được: $24 + 36 = 60$ gói bánh.
\textbf{Đáp án của câu hỏi là B}
}
\end{ex}

\begin{ex} %[3X2V28]
An có 15 viên bi. Bình có số bi gấp đôi số bi của An. Cường có số bi ít hơn Bình 8 viên. Hỏi Cường có bao nhiêu viên bi?
\choice
{$30$ viên bi}
{$25$ viên bi}
{$18$ viên bi}
{\True $22$ viên bi}
\loigiai{
Bình có: $15 \times 2 = 30$ viên bi.\
Cường có: $30 - 8 = 22$ viên bi.
\textbf{Đáp án của câu hỏi là D}
}
\end{ex}

\begin{ex} %[3X2V28]
Mẹ mua một con gà nặng $2\text{ kg}$ và một con vịt nặng gấp 3 lần con gà. Hỏi cả gà và vịt nặng tất cả bao nhiêu ki-lô-gam?
\choice
{\True $8\text{ kg}$}
{$6\text{ kg}$}
{$7\text{ kg}$}
{$9\text{ kg}$}
\loigiai{
Con vịt nặng: $2 \times 3 = 6\text{ (kg)}$. Cả hai con nặng: $2 + 6 = 8\text{ (kg)}$.
\textbf{Đáp án của câu hỏi là A}
}
\end{ex}

\begin{ex} %[3X2V28]
Một thùng chứa $45\ell$ dầu. Người ta đã lấy ra 3 can dầu, mỗi can chứa $5\ell$ dầu. Hỏi trong thùng còn lại bao nhiêu lít dầu?
\choice
{$35\ell$}
{$25\ell$}
{\True $30\ell$}
{$40\ell$}
\loigiai{
Số dầu đã lấy ra là: $5 \times 3 = 15\text{ (\ell)}$.\
Số dầu còn lại trong thùng là: $45 - 15 = 30\text{ (\ell)}$.
\textbf{Đáp án của câu hỏi là C}
}
\end{ex}

\begin{ex} %[3X2V28]
Mai có 36 bông hoa. Mai tặng Lan $\dfrac{1}{4}$ số hoa đó rồi tặng Cúc thêm 5 bông hoa nữa. Mai còn lại bao nhiêu bông hoa?
\choice
{$27$ bông hoa}
{\True $22$ bông hoa}
{$24$ bông hoa}
{$20$ bông hoa}
\loigiai{
Mai tặng Lan: $36 : 4 = 9$ bông hoa.\
Mai còn lại sau khi tặng cả hai bạn: $36 - (9 + 5) = 36 - 14 = 22$ bông hoa.
\textbf{Đáp án của câu hỏi là B}
}
\end{ex}

\begin{ex} %[3X2V28]
Lớp 3A có 18 bạn nam. Số bạn nữ ít hơn số bạn nam là 4 bạn. Hỏi lớp 3A có tất cả bao nhiêu học sinh?
\choice
{$28$ bạn}
{$30$ bạn}
{$34$ bạn}
{\True $32$ bạn}
\loigiai{
Số bạn nữ là: $18 - 4 = 14$ bạn. Tổng số học sinh cả lớp là: $18 + 14 = 32$ bạn.
\textbf{Đáp án của câu hỏi là D}
}
\end{ex}

\begin{ex} %[3X2V28]
Bác thợ mộc cưa một cây gỗ dài $48\text{ dm}$ thành các đoạn dài $8\text{ dm}$. Mỗi lần cưa mất 5 phút. Hỏi bác cưa xong cây gỗ hết tất cả bao nhiêu phút?
\choice
{\True $25$ phút}
{$30$ phút}
{$20$ phút}
{$35$ phút}
\loigiai{
Cây gỗ được cưa thành: $48 : 8 = 6$ đoạn. Bác cần thực hiện: $6 - 1 = 5$ lần cưa.\
Tổng thời gian cưa là: $5 \times 5 = 25$ phút.
\textbf{Đáp án của câu hỏi là A}
}
\end{ex}

\begin{ex} %[3X2V28]
Một người nuôi 14 con gà mái. Số gà trống nhiều hơn gà mái là 6 con. Người đó bán đi một nửa tổng số gà. Hỏi người đó đã bán bao nhiêu con gà?
\choice
{$20$ con gà}
{$34$ con gà}
{\True $17$ con gà}
{$15$ con gà}
\loigiai{
Số gà trống là: $14 + 6 = 20$ con. Tổng số gà: $14 + 20 = 34$ con.\
Số gà đã bán là: $34 : 2 = 17$ con gà.
\textbf{Đáp án của câu hỏi là C}
}
\end{ex}

\begin{ex} %[3X2V28]
Có 4 xe chở học sinh đi tham quan, mỗi xe chở 24 bạn. Khi đến nơi có thêm 12 bạn đi xe ô tô con đến nhập đoàn. Hỏi có tất cả bao nhiêu bạn đi tham quan?
\choice
{$96$ bạn}
{\True $108$ bạn}
{$100$ bạn}
{$112$ bạn}
\loigiai{
Số bạn trên 4 xe khách: $24 \times 4 = 96$ bạn. Tổng số bạn tham quan: $96 + 12 = 108$ bạn.
\textbf{Đáp án của câu hỏi là B}
}
\end{ex}

\begin{ex} %[3X2V28]
Đoạn đường từ nhà đến công viên dài $60\text{ m}$. Bạn Hùng đi bộ được $\dfrac{1}{3}$ quãng đường thì gặp bạn Cường, hai bạn cùng đi bộ thêm $15\text{ m}$ nữa thì dừng lại nghỉ. Hỏi hai bạn còn cách công viên bao nhiêu mét?
\choice
{$20\text{ m}$}
{$35\text{ m}$}
{$30\text{ m}$}
{\True $25\text{ m}$}
\loigiai{
Quãng đường Hùng đã đi trước khi gặp Cường: $60 : 3 = 20\text{ m}$. Tổng quãng đường đã đi: $20 + 15 = 35\text{ m}$. Khoảng cách còn lại: $60 - 35 = 25\text{ m}$.
\textbf{Đáp án của câu hỏi là D}
}
\end{ex}

\begin{ex} %[3X2V28]
Một người mua 5 quyển vở cùng loại hết 45 nghìn đồng. Người đó mua thêm 3 quyển vở như thế nữa. Tổng số tiền người đó phải trả cho tất cả 8 quyển vở là
\choice
{\True $72$ nghìn đồng}
{$65$ nghìn đồng}
{$80$ nghìn đồng}
{$60$ nghìn đồng}
\loigiai{
Giá tiền một quyển vở là: $45 : 5 = 9$ nghìn đồng.\
Tổng số tiền mua 8 quyển vở là: $9 \times 8 = 72$ nghìn đồng.
\textbf{Đáp án của câu hỏi là A}
}
\end{ex}

\begin{ex} %[3X2V28]
Thùng to có $56\ell$ nước, thùng bé có $24\ell$ nước. Hỏi phải chuyển từ thùng to sang thùng bé bao nhiêu lít nước để hai thùng có lượng nước bằng nhau?
\choice
{$32\ell$}
{$8\ell$}
{\True $16\ell$}
{$20\ell$}
\loigiai{
Thùng to hơn thùng bé: $56 - 24 = 32\text{ (\ell)}$. Cần chuyển từ thùng to sang thùng bé: $32 : 2 = 16\text{ (\ell)}$.
\textbf{Đáp án của câu hỏi là C}
}
\end{ex}

\textbf{Phần 2. Trắc nghiệm đúng sai}
\setcounterpageref{ex}{0}

\begin{ex} %[3X2V28]
Xét tính đúng sai của các bước giải bài toán bằng hai bước tính:
\choiceTF
{\True Bài toán giải bằng hai bước tính luôn yêu cầu tìm một đại lượng trung gian trước khi tìm đáp số cuối cùng.}
{Chỉ có thể kết hợp phép cộng với phép trừ, không thể kết hợp phép nhân với phép chia.}
{\True Cần đọc kĩ đề bài để xác định đại lượng đã cho và đại lượng cần tìm.}
{Hai bước tính luôn cho ra hai kết quả giống hệt nhau.}
\loigiai{
\begin{itemize}
    \item Ý a) Đúng: Đặc điểm phương pháp giải.
    \item Ý b) Sai: Sai vì có thể kết hợp bất kì phép tính nào.
    \item Ý c) Đúng: Kĩ năng giải toán cơ bản.
    \item Ý d) Sai: Sai vì mỗi bước tính tìm một đối tượng khác nhau.
\end{itemize}
\begin{itemize}
    \item a) \textbf{Đúng}.
    \item b) \textbf{Sai}.
    \item c) \textbf{Đúng}.
    \item d) \textbf{Sai}.
\end{itemize}
}
\end{ex}

\begin{ex} %[3X2V28]
Nhà Lan thu hoạch được 3 bao khoai tây, mỗi bao nặng $18\text{ kg}$. Nhà Lan giữ lại $14\text{ kg}$ để ăn, số còn lại đem bán:
\choiceTF
{\True Tổng khối lượng khoai tây thu hoạch được là $54\text{ kg}$.}
{\True Số ki-lô-gam khoai tây đem bán là $40\text{ kg}$.}
{Số khoai tây đem bán ít hơn số khoai tây giữ lại để ăn.}
{\True Nếu mỗi ki-lô-gam khoai tây bán được với giá nhất định thì tiền thu được từ 40 kg khoai tây.}
\loigiai{
\begin{itemize}
    \item Ý a) Đúng: $18 \times 3 = 54\text{ kg}$.
    \item Ý b) Đúng: $54 - 14 = 40\text{ kg}$.
    \item Ý c) Sai: Sai vì 40 kg nhiều hơn 14 kg.
    \item Ý d) Đúng: Đúng với số khoai tây bán.
\end{itemize}
\begin{itemize}
    \item a) \textbf{Đúng}.
    \item b) \textbf{Đúng}.
    \item c) \textbf{Sai}.
    \item d) \textbf{Đúng}.
\end{itemize}
}
\end{ex}

\textbf{Phần 3. Trả lời ngắn}
\setcounterpageref{ex}{0}

\begin{ex} %[3X2V28]
Cửa hàng có 3 thùng dầu, mỗi thùng chứa $25\ell$. Cửa hàng đã bán được $38\ell$. Hỏi cửa hàng còn lại bao nhiêu lít dầu?
\shortans{37}
\loigiai{
\textbf{Đáp số:} $37$\\
Tổng số dầu: $25 \times 3 = 75\ell$. Còn lại: $75 - 38 = 37\text{ (\ell)}$.
}
\end{ex}

\begin{ex} %[3X2V28]
Nam có 16 viên bi. Dũng có số bi gấp 3 lần Nam. Hỏi cả hai bạn có tất cả bao nhiêu viên bi?
\shortans{64}
\loigiai{
\textbf{Đáp số:} $64$\\
Dũng có: $16 \times 3 = 48$ viên. Cả hai bạn có: $16 + 48 = 64$ viên bi.
}
\end{ex}

\begin{ex} %[3X2V28]
Mẹ mua $4\text{ kg}$ táo hết 80 nghìn đồng. Hỏi mẹ mua $3\text{ kg}$ táo như thế hết bao nhiêu nghìn đồng?
\shortans{60}
\loigiai{
\textbf{Đáp số:} $60$\\
Một kg táo giá: $80 : 4 = 20$ nghìn đồng. Ba kg giá: $20 \times 3 = 60$ nghìn đồng.
}
\end{ex}

\begin{ex} %[3X2V28]
Lớp 3A trồng được 32 cây xanh, lớp 3B trồng được nhiều hơn lớp 3A là 8 cây xanh. Hỏi cả hai lớp trồng được tất cả bao nhiêu cây xanh?
\shortans{72}
\loigiai{
\textbf{Đáp số:} $72$\\
Lớp 3B trồng: $32 + 8 = 40$ cây. Cả hai lớp trồng: $32 + 40 = 72$ cây xanh.
}
\end{ex}

\begin{ex} %[3X2V28]
Một tấm vải dài $60\text{ m}$. Người ta cắt ra $\dfrac{1}{3}$ tấm vải để may áo, sau đó cắt tiếp $15\text{ m}$ để may rèm. Hỏi tấm vải còn lại bao nhiêu mét?
\shortans{25}
\loigiai{
\textbf{Đáp số:} $25$\\
May áo hết: $60 : 3 = 20\text{ m}$. Tấm vải còn lại: $60 - (20 + 15) = 25\text{ (m)}$.
}
\end{ex}

\begin{ex} %[3X2V28]
Bác Năm nuôi 2 chuồng thỏ, mỗi chuồng có 12 con. Bác mua thêm 8 con thỏ nữa rồi chia đều tất cả số thỏ vào 4 lồng mới. Hỏi mỗi lồng mới có bao nhiêu con thỏ?
\shortans{8}
\loigiai{
\textbf{Đáp số:} $8$\\
Tổng số thỏ: $12 \times 2 + 8 = 32$ con. Mỗi lồng có: $32 : 4 = 8$ con thỏ.
}
\end{ex}
